\originalchapter{期中测验}\label{ch:midterm}
原件3页, 题面2页, 5题共40分. 原开放书籍及禁止协作规定见课程信息; 以下是归档学习用的 AI 独立解答.
\begin{exercise}\label{midterm:1}
原题1（8分）. 对 $a<b$, 令 $\Lambda([a,b])$ 为所有 Lipschitz 连续函数, 即存在 $L\ge0$ 使 $|f(x)-f(y)|\le L|x-y|$. 证明它在范数 $\norm f=\norm f_\infty+\sup_{x\ne y}|f(x)-f(y)|/|x-y|$ 下为 Banach 空间. 原题允许使用上确界范数下 $C([a,b])$ 完备.
\end{exercise}
\begin{proof}[AI 独立解答]
记第二项为 $[f]_{\rm Lip}$. 和及数乘的 Lipschitz 常数受相应常数的和及绝对数乘控制, 所以这是向量空间. 上确界项保证正定性, 两项分别齐次且满足三角不等式, 因而确为范数.

设 $f_n$ 为本范数 Cauchy 列. 它在上确界范数下收敛于连续 $f$. 存在 $M$ 使 $[f_n]_{\rm Lip}\le M$, 对每对 $x,y$ 取极限得 $|f(x)-f(y)|\le M|x-y|$, 所以 $f\in\Lambda$. 又对 $n,m\ge N$, 有 $[f_n-f_m]_{\rm Lip}<\eps$. 固定 $n,x,y$, 令 $m\to\infty$, 再对 $x\ne y$ 取上确界得 $[f_n-f]_{\rm Lip}\le\eps$. 与一致收敛结合, 得本范数收敛.
\end{proof}
\begin{exercise}\label{midterm:2}
原题2（8分）. 对 $a\in\ell^\infty$, 证明 $\norm{a+c_0}_{\ell^\infty/c_0}=\limsup_{n\to\infty}|a_n|=\inf_{n\ge1}\sup_{k\ge n}|a_k|$. 原题允许使用: 加上趋零数列不改变绝对值的上极限.
\end{exercise}
\begin{proof}[AI 独立解答]
取 $b^{(N)}=(-a_1,\ldots,-a_N,0,\ldots)\in c_0$, 则商范数至多 $\sup_{k>N}|a_k|$, 令 $N\to\infty$ 得上界. 对任意 $b\in c_0$, $\limsup|a_n|=\limsup|a_n+b_n|\le\norm{a+b}_\infty$; 对 $b$ 取下确界得下界. 许可使用的上极限恒等式也由 $\big||a_n+b_n|-|a_n|\big|\le|b_n|\to0$ 直接得到.
\end{proof}
\begin{exercise}\label{midterm:3}
原题3. (a)（4分）若 $V,W$ 为 Banach 空间, $T_n\in\B(V,W)$, 线性 $T:V\to W$ 满足 $T_nx\to Tx$ 对每个 $x\in V$ 成立, 证明 $T$ 有界. 不假设算子范数收敛. (b)（4分）若同一向量空间的两范数满足 $\norm v_1\le\norm v_2$, 且两范数下都完备, 证明存在 $C>0$ 使 $\norm v_2\le C\norm v_1$.
\end{exercise}
\begin{proof}[AI 独立解答]
(a) 每条收敛的 $(T_nx)$ 有界. 统一有界原理\ref{thm:ub}给出 $M=\sup_n\norm{T_n}<\infty$, 故 $\norm{Tx}=\lim_n\norm{T_nx}\le M\norm x$. 此复用推论\ref{cor:stronglimit}.

(b) 恒等映射 $I:(V,\norm\cdot_2)\to(V,\norm\cdot_1)$ 为有界双射, 两端完备. 有界逆定理给出其逆有界, 即所需不等式; 零空间任选 $C>0$. 这也复用第\ref{ch:baire}章两完备范数的习题解答.
\end{proof}
\begin{exercise}\label{midterm:4}
原题4. 设 $E\subset\R$ 可测, $f,g:E\to\R$ 可测, 定义 $\norm f_\infty=\inf\{C\ge0:m(\{|f|>C\})=0\}$, 空集下确界取 $+\infty$. (a)（4分）证明 $|f|\le\norm f_\infty$ 几乎处处. (b)（4分）对 $c\in\R$, 证明 $\norm{cf}_\infty=|c|\norm f_\infty$, $\norm{f+g}_\infty\le\norm f_\infty+\norm g_\infty$.
\end{exercise}
\begin{proof}[AI 独立解答]
(a) 若本质上确界无穷, 不等式处处成立. 若 $M<\infty$, 对每个 $n$ 选可行本质界小于 $M+1/n$, 除去其零测例外的可数并, 得 $|f|\le M+1/n$ 对所有 $n$ 成立, 令 $n\to\infty$.

(b) $c\ne0$ 时可行界经乘 $|c|$ 一一对应, 包括没有有限可行界的情形. $c=0$ 时 $cf=0$ 的范数为零, 等式使用非负扩展实数约定 $0\cdot\infty=0$. 对三角不等式, 任一右端项为无穷时自动成立; 否则在两个本质界的共同零测例外外用逐点三角不等式, 再取本质界下确界.
\end{proof}
\begin{exercise}\label{midterm:5}
原题5. 设非负简单函数 $\phi=\sum_{i=1}^n a_i\ind_{A_i}$, $A_i\subset E$ 为可测分割. (a)（4分）若 $F\subset E$ 可测, 证明 $\int_F\phi=\int_E\phi\ind_F$. (b)（4分）若 $F,G\subset E$ 可测且不交, 证明 $\int_{F\cup G}\phi=\int_F\phi+\int_G\phi$.
\end{exercise}
\begin{proof}[AI 独立解答]
(a) $\ind_{A_i\cap F}=\ind_{A_i}\ind_F$, 所以两积分都等于 $\sum_i a_i m(A_i\cap F)$. (b) $A_i\cap(F\cup G)$ 为 $A_i\cap F$ 与 $A_i\cap G$ 的不交并; 对测度用可加性, 再作有限个非负求和, 即得等式. 所有加法允许无穷, 零系数乘无穷按定义取零.
\end{proof}
